% ============================================================================
\chapter{用户态 libc}
\label{ch:libc}
% Stage F-8 ~ F-9
% ============================================================================

\section{本章目标}

\begin{itemize}
  \item 理解 libc 在用户程序与内核之间的桥梁角色
  \item 对比 musl / glibc / newlib，选择合适的策略
  \item 实现 C 运行时启动代码（\texttt{crt0.o}、\texttt{\_start} → \texttt{main}）
  \item 实现精简 libc 核心函数（printf/malloc/fopen/string.h）
  \item 验证用户态程序链接 libc 并正常运行
\end{itemize}

% ============================================================================
\section{概念：libc 的角色}
\label{sec:libc-concept}
% ============================================================================

% TODO:
% - libc = 用户态对内核 syscall 的封装 + C 标准库函数
% - 分层：用户程序 → libc API（printf, malloc） → syscall wrapper → 内核
% - 没有 libc 时用户程序只能 inline asm 调 syscall → 极不方便
% - libc 提供的核心功能：
%     I/O：printf/scanf/fopen/fread/fwrite/fclose
%     内存：malloc/free/calloc/realloc（底层调 brk/mmap）
%     字符串：strlen/strcpy/strcmp/memcpy/memset
%     进程：exit/fork/exec/wait（syscall 的 thin wrapper）
% - 方案对比：
%     musl：轻量、代码清晰、适合嵌入式和自制 OS
%     glibc：功能齐全但庞大、不适合早期移植
%     newlib：专为嵌入式设计，需实现少量 syscall stub
%     自写精简 libc：最小依赖，全部可控（我们先用这个）

\subsection{C 运行时启动流程}

% TODO:
% - 用户程序入口不是 main()，而是 _start（crt0.o 提供）
% - _start → __libc_init（初始化 heap/stdio/env）→ main(argc, argv, envp) → exit(ret)
% - crt0.o 汇编：
%     _start:
%         xor ebp, ebp       ; 栈底标记（调试器用）
%         pop ecx             ; argc（内核在栈上准备的）
%         mov esi, esp        ; argv 起始地址
%         call main
%         push eax            ; main 返回值
%         call exit           ; 不返回

% ============================================================================
\section{syscall wrapper 层}
\label{sec:libc-syscall}
% ============================================================================

% TODO:
% - 每个 POSIX 函数的 thin wrapper：
%     ssize_t write(int fd, const void* buf, size_t count) {
%         return syscall3(SYS_write, fd, (uint32_t)buf, count);
%     }
% - 错误处理：syscall 返回负数 → 设 errno = -ret, 返回 -1
% - errno 定义（ENOSYS, EBADF, ENOMEM, EFAULT, ...）

% ============================================================================
\section{核心函数实现}
\label{sec:libc-impl}
% ============================================================================

% TODO:
% - stdio：
%     printf → vsnprintf + write(1, ...)
%     fopen → open + 分配 FILE 结构
%     fread/fwrite → read/write + 缓冲
% - malloc/free：
%     小块 → brk 扩展堆 + 简单 free-list
%     大块 → mmap(MAP_ANONYMOUS)
% - string.h：strlen, strcpy, strcmp, memcpy, memset, strncmp, ...
% - ctype.h：isalpha, isdigit, isspace, toupper, tolower
% - 静态链接：ar rcs libc.a *.o → 用户程序 -lc 链接

% ============================================================================
\section{验证}
\label{sec:libc-verify}
% ============================================================================

% TODO:
% - 测试 1：printf("hello %d %s\n", 42, "world") → 串口输出正确
% - 测试 2：malloc(1024) → memset → free → malloc → 地址复用
% - 测试 3：fopen("/etc/hello") → fread → fclose → 内容匹配
% - 测试 4：main 返回 0 → exit(0) 正常退出

% ============================================================================
\section{本章小结}
% ============================================================================

精简 libc 为用户程序提供了标准 C 编程接口，屏蔽了 syscall 的底层细节。
crt0.o + syscall wrapper + 核心函数三层结构让用户程序能用熟悉的 printf/malloc 编程。
下一章实现用户态 Shell，替代内核内置的 shell，迈向真正的多用户系统。
